\scanpage{249}
\mspage{\unsure{195|175}}

se rendre compte à présent que $\unsure{\Phi|\mathscr{D}}_\infty$ a un avantage exactement du $=$ ordre sur $C$, en introduisant
\[
\begin{aligned}
W_{\unsure{\omega|w}}''={}&\text{plus petite partie }\Sigma\text{ de }\mathrm{Fl}(\mathrm{Cat})\\
&\text{satisfaisant }W(1,2\ \mathrm{bis},3,6\ \mathrm{bis})
\end{aligned}
\]
et noter que
\[
\Phi_\infty(X)\in\mathrm{Ob}\ \underline{\mathrm{Perf}}_{W_\omega''}S,\qquad
\left|\begin{aligned}
i_X&\in\bigl(W_\omega''\deleted{\unsure{\bigr)_S^f}}\bigr.\\[-2pt]
&\deleted{\unsure{(W_\omega'')_S^f}}
\end{aligned}\right.
\]
-- mais là je m'aperçois qu'on a
\[
W_\omega^{\unsure{\prime\prime\prime|\prime|b}}\subset W_\omega''
\]
donc $C$ est quand même un peu meilleur à cet égard \ldots

Quand on dispose de $\Phi_\infty$ avec les propriétés qu'on vient d'expliciter, on trouve que \deleted{\unsure{les}} l'inclusion
\[
\underline{\mathrm{Perf}}_W S\to\mathrm{Fib}_W(S)
\]
% slanted
\margin{On peut y remplacer $=$ $\underline{\mathrm{Perf}}_W S$ par $\underline{\mathrm{Perf}}_{W_\omega''}S$, \uline{ou} par toute \inserted{autre} sous-cat. pleine de $\underline{\mathrm{Fib}}_W(S)$ contenant $\underline{\mathrm{Perf}}_{W_\omega''}S$.}
induit une équivalence non seulement pour la localisation par $W_S$, mais aussi pour les localisés fins
\[
\begin{aligned}
(\underline{\mathrm{Perf}}_W S)(W_S^f)^{-1}&\longrightarrow(\mathrm{Fib}_W(S))(W_S^f)^{-1}.\\
&\qquad\overset{\mathrm{def}}{=}\mathrm{HOT}_W^{\unsure{\mathrm{lc}|\mathrm{le}|\mathrm{loc}}}(S)
\end{aligned}
\]

Mais on ne peut \deleted{\unsure{pour}} autant réduire le nombre des \unsure{localisés|localisions} \unsure{utiles} à trois, en l'absence de $W(4\mathrm b)$. Tout au plus, ayant un foncteur $\Phi_\infty$ qui est plus souple que $C$, en l'absence de $W(4\mathrm b)$, du fait qu'on sait de toute façon que
\[
i_X:X\longrightarrow\Phi_\infty(X)
\]
est dans $W_S^f$ si $X$ est dans $\underline{\mathrm{Fib}}_W S$, peut-on
% archive sheet number: 248
